\section{Discussion}
\label{sec:discussion}

\paragraph{Synthesis.}
Our experiments reveal a coherent picture of how action-relevant structure flows through a VLA.
The vision encoder is the \emph{source}: VL-pretrained encoders provide geometric structure aligned with the action space (SigLIP+Projector RSA $\rho = 0.615$; CLIP $\rho = 0.557$), far exceeding what the LLM produces at any layer ($\rho \leq 0.501$).
The LLM is a \emph{router}: it degrades distributed geometric structure ($\rho$: $0.615 \to 0.501$; $R^2$: $0.706 \to 0.683$) while concentrating a functionally superior representation at the readout token ($R^2 = 0.731$).
Training teaches the \emph{routing policy}: the geometric structure itself is invariant to training ($\Delta\rho = 0.021$, $p = 0.20$), but functional concentration at the readout requires it ($\Delta R^2 = 0.045$, $p < 0.0001$).
Notably, the projector alone captures 96.6\% of the trained readout's functional capacity ($0.706/0.731$), suggesting that the LLM's contribution, while statistically significant, is incremental.

\paragraph{VL pretraining as a general source of action-relevant geometry.}
The inclusion of CLIP ViT-L/14 as an external baseline provides evidence that action-relevant geometric structure is a general product of vision-language pretraining, not an artifact of SigLIP's specific architecture or training data.
Both VL-pretrained encoders---SigLIP ($\rho = 0.615$) and CLIP ($\rho = 0.557$)---substantially exceed vision-only DINOv2 ($\rho = 0.388$).
The 0.169-point gap between CLIP and DINOv2, both ViT-L/14 architectures differing primarily in pretraining objective, isolates language grounding as the operative factor.
The additional increment from CLIP to SigLIP+Projector ($\Delta = 0.058$) is consistent with the projector's VLA-supervised alignment providing a modest boost, but the dominant contribution ($\rho = 0.557 \gg 0.388$) originates in VL pretraining itself.

\paragraph{Geometric-functional dissociation.}
A recurring finding across our analyses is the dissociation between geometric alignment (RSA) and functional decodability (probe $R^2$).
DINOv2 achieves probe performance statistically indistinguishable from the SigLIP+Projector ($R^2$: $0.712$ vs.\ $0.706$, $p = 0.22$) despite substantially lower RSA ($0.388$ vs.\ $0.615$).
This dissociation implies that action information---in the sense of decodable content---is broadly available in strong visual features regardless of pretraining objective; what VL pretraining adds is \emph{geometric alignment}, organizing that information into a structure that mirrors the action space.
For VLA design, this distinction matters: encoder selection determines not just what information is available, but how it is organized, with downstream consequences for how easily the LLM can route it.

\paragraph{Metric sensitivity: RSA versus CKA.}
The divergence between RSA and linear CKA underscores the importance of multi-metric evaluation.
Linear CKA produces a ranking---untrained LLM ($0.089$) $>$ trained ($0.083$) $>$ SigLIP+Projector ($0.081$) $>$ DINOv2 ($0.066$)---that partially reverses the RSA ordering.
Under linear CKA, training appears to \emph{reduce} similarity to the action space, an interpretation that contradicts both RSA and functional probe evidence.
RDM kernel CKA resolves this discrepancy, restoring the RSA-consistent ranking (trained $0.260$ $>$ Proj $0.257$ $>$ untrained $0.246$ $>$ DINOv2 $0.241$).
The divergence likely reflects linear CKA's sensitivity to second-order statistics and its known vulnerability to dimension mismatch between high-dimensional representations and the 7-D action space~\cite{kornblith2019}.
The practical implication is that the LLM transforms rather than destroys action structure, but the narrative one draws depends critically on which metric is employed.

\paragraph{Implications for VLA design.}
Our findings suggest four practical considerations.
\emph{First}, vision encoder pretraining quality is the primary lever for action-relevant representations; investing in stronger VL encoders may yield larger gains than scaling the LLM backbone.
\emph{Second}, for downstream modules that operate on distributed representations (e.g., auxiliary losses, multi-task heads), the projector output may be preferable to LLM-processed features, given its higher geometric fidelity and 96.6\% of readout-level decodability.
\emph{Third}, the LLM's value lies in readout concentration---routing distributed information to a single decoding site---which informs action head design: architectures that exploit this concentration (e.g., reading from a dedicated token rather than pooling LLM features) may better leverage the LLM's trained routing.
\emph{Fourth}, the invariance of representation geometry to language instructions ($\Delta\rho = 0.004$, n.s.; content controls within $\pm 0.017$) suggests that instructions influence action selection through a mechanism other than modulating visual representation structure, inviting further investigation into how language conditions behavior in VLAs.

\paragraph{Connection to behavioral evidence.}
Our representational findings complement the behavioral results of \citet{grant2026}, who showed that VLA policies are disproportionately driven by visual input and that language instructions are ``encoded but not used'' for action selection.
Our data provide a representational account of these behavioral patterns: vision dominates because the frozen encoder supplies the action-relevant geometric structure that the LLM routes to the readout, while instruction content does not modulate this geometric structure (all four instruction conditions yield RSA within $\pm 0.017$).
The alignment between representational and behavioral evidence strengthens the case that VL-pretrained vision encoders, rather than LLM reasoning, are the primary computational substrate for action generation in current VLA architectures.

\paragraph{Limitations.}
Several limitations scope our conclusions.
We study a single VLA architecture (OpenVLA-7B) on a single benchmark suite (LIBERO); generalization to other architectures (e.g., RT-2~\cite{brohan2023_rt2}) and domains (e.g., mobile manipulation, dexterous hands) remains to be established.
RSA captures structural alignment between representations and actions but does not establish causal contribution---high RSA does not guarantee that the representation is causally used by the policy.
The CLIP and DINOv2 action probes were trained with different hyperparameters (100 epochs) than the projector and LLM probes (50 epochs); while within-experiment comparisons are valid, absolute $R^2$ values are not directly comparable across these two probe sets, and we have been careful to restrict cross-experiment claims accordingly.
Finally, our analyses characterize frozen representations post hoc; they do not address how these representations emerge during VLA pretraining or fine-tuning, nor whether architectural interventions (e.g., unfreezing the vision encoder) would alter the source--router dynamic we observe.
